\section{Normed Vector Spaces}

  Given a vector space $V$, we can induce different structures on it to allow us to conduct different measurements on it. For example, the endowment of the basis on $V$ allows us to represent vector as an $n$-tuple of scalars, and from topology, we can also define an arbitrary topology or metric on a vector space as well. 

  \begin{definition}[Norm]
    A \textbf{norm} on a vector space $V$ over field $\mathbb{C}$ is a mapping 
    \begin{equation}
      \rho: V \longrightarrow \mathbb{R}
    \end{equation}
    satisfying three properties 
    \begin{enumerate}
      \item $\rho (x) \geq 0$, with $\rho(x) = 0 \iff x = 0$
      \item For $a \in \mathbb{C}$, $\rho (a x) = |a| \rho(x)$ 
      \item $\rho(x + y) \leq \rho(x) + \rho(y)$ 
    \end{enumerate}
    A norm allows us to define some notion of a magnitude or length on each vector in $V$. A vector space $V$ with a norm is called a \textbf{normed space}, denoted $(V, \rho)$. 
  \end{definition}

  \begin{example}[Absolute Value]
    The absolute value function $|\cdot|: \mathbb{C} \longrightarrow \mathbb{R}_+$ is an example of a norm on the 1 dimensional space $\mathbb{C}$. 
  \end{example}

  \begin{example}[Euclidean Norm, $L_2$-Norm]
    The Euclidean norm of a vector $x \equiv (x_1, x_2, \cdots, x_n)^T \in \mathbb{R}^n$ is defined
    \begin{equation}
      ||x||_2 \equiv \bigg( \sum_{i=1}^n x_i^2 \bigg)^{\frac{1}{2}}
    \end{equation}
    This is the most commonly used norm in $\mathbb{R}^n$. 
  \end{example}

  \begin{example}[Taxicab Norm, Manhattan Norm]
    The Taxicab norm of $x \equiv (x_1, x_2, \cdots, x_n)^T \in \mathbb{R}^n$ is defined
    \begin{equation}
      ||x||_1 \equiv \sum_{i=1}^n |x_i|
    \end{equation}
  \end{example}

  \begin{example}[Infinity Norm, $L_\infty$-Norm]
    The Infinity norm of vector $x \equiv (x_1, x_2, \cdots, x_n)^T \in \mathbb{R}^n$ is defined
    \begin{equation}
      ||x||_\infty \equiv \max{\{|x_1|, |x_2|, \cdots, |x_n|\}}
    \end{equation}
  \end{example}

  \begin{example}[p-norm, $L_p$-Norm]
    Let $p\geq 1$ be a real number. The p-norm of a vector 
    \begin{equation}
      x \equiv (x_1, x_2, \cdots, x_n)^T \in \mathbb{R}^n
    \end{equation}
    is defined
    \begin{equation}
      ||x||_p \equiv \bigg( \sum_{i=1}^n x_i^p \bigg)^{\frac{1}{p}}
    \end{equation}
    For $0<p<1$, this function could be of some use, but it is not considered a norm since it violates the triangle inequality. When $p = 1$ and $p =2$, the norm is the Taxicab norm and Euclidean norm, respectively, and 
    \begin{equation}
      \lim_{p \rightarrow \infty} ||\cdot||_p = ||\cdot||_\infty
    \end{equation}
  \end{example}

  \begin{theorem}[Norm Induces Metric]
    Every norm induces a metric in the following way
    \begin{equation}
      d(x, y) \equiv \rho(x-y)
    \end{equation}
    However, a metric does not necessarily induce a norm because the definition
    \begin{equation}
      \rho(x) \equiv d(x, 0)
    \end{equation}
    is not guaranteed to have all properties of the norm. 
  \end{theorem}

\subsection{Dual Spaces}

\subsection{Norms of Linear Mappings}

  Since the algebra of linear operators is itself a vector space, we can also define structures on it, too. We focus on matrix norms. 

  \begin{definition}[Operator Norm]
    Let $A: X \longrightarrow U$ be linear. Then, we define
    \begin{equation}
      \|A\| = \sup_{\|x\|=1} \|A x\|
    \end{equation}
    Note that $\|A x\|$ is measure with respect to the norm of $U$ and $\|x\|$ the norm of $X$. 
  \end{definition}

  There is a very nice visualization of this. 

  \begin{figure}[H]
    \centering 
    \begin{tikzpicture}
      % Left diagram (Unit circle in R^n)
      \begin{scope}[shift={(-4,0)}]
        % Axes
        \draw[-] (-2.5,0) -- (2.5,0);
        \draw[-] (0,-2.5) -- (0,2.5);
        
        % Unit circle
        \draw[thick] (0,0) circle (1.5cm);
        
        % Label R^n
        \node at (0.8,2) {$\mathbb{R}^n$};
      \end{scope}
      
      % Arrow with label A
      \draw[-{Stealth[length=3mm,width=2mm]}] (-1.5,0.5) -- (1.5,0.5);
      \node at (0,1) {$A$};
      
      % Right diagram (Ellipse in R^m)
      \begin{scope}[shift={(4,0)}]
        % Axes
        \draw[-] (-2.5,0) -- (2.5,0);
        \draw[-] (0,-2.5) -- (0,2.5);
        
        % Rotated Ellipse (30 degrees)
        \draw[thick, rotate=30] (0,0) ellipse (2cm and 1cm);
        
        % Major axis line
        \draw[->, thick, rotate=30] (0,0) -- (2,0);
        
        % Label R^m
        \node at (1.8,2) {$\mathbb{R}^m$};
      \end{scope}
    \end{tikzpicture}
    \caption{The norm of $A$ is the length of the major axis of the ellipsoid. Given that $\dim{X}=n$, imagine the $n$-dimensional unit ball in $X$ being transformed under $A$. The image of the ball should be an ellipsoid (of dimension $\leq m$) in $U$. } 
    \label{fig:operator_norm}
  \end{figure}

  \begin{theorem}
    \begin{align}
      \|A z\| \leq \|A\| \|z\| \text{ for all } z \in X \\
      \|A\| = \sup_{\|x\|, \|v\| = 1} (A x, v)
    \end{align}
  \end{theorem}
  \begin{proof}
    \begin{align*}
      &\|A z\| \leq \sup{\|A z\|} = \sup{\Big|\Big| A \frac{z}{\|z\|} \Big|\Big|} = \|A\| \|z\| \\
      &\|u\| \equiv \max_{\|v\|=1} (u, v) \implies \|A x\| \equiv \max_{\|v\|=1} (Ax, v) \implies \|A\| \equiv \sup_{\|x\|, \|v\| =1} (A x, v)
    \end{align*}
  \end{proof}

  \begin{theorem}[Properties of Matrix Norm]
    Let there exist any $k \in \mathbb{F}$, with any $A, B: X \longrightarrow U$, $C: U \longrightarrow V$. Then, 
    \begin{enumerate}
      \item $\|k A\| = |k| \|A\|$
      \item $\|A + B\| \leq \|A\| + \|B\|$
      \item $\|C A\| \leq \|C\| \|A\|$
      \item $\|A\| = \|A^\dagger\|$
    \end{enumerate}
  \end{theorem}

  \begin{theorem}
    A simple lower and upper bound of $\|A\|$ can be defined
    \begin{equation}
      r(A) \leq \|A\| \leq \bigg( \sum_{i, j} a_{i j}^2 \bigg)^\frac{1}{2}
    \end{equation}
  \end{theorem}

  Matrix norms have extremely useful applications in determining the existence of inverses. 

  \begin{theorem}
    Let $A$ be invertible and 
    \begin{equation}
      \|A - B\| < \frac{1}{\|A^{-1}\|}
    \end{equation}
    in the sense that $B$ is "close" to $A$. Then $B$ is invertible. 
  \end{theorem}


